% !PRINTED 609
\weberpage{609}
Ist nun
\weq{15}{\alpha_1 = a_{1,1}, \quad \alpha_2 = a_{1,2}+a_{2,2}\Theta}
eine Basis von $\varphi$, so ist andererseits (nach \S.\,146)
\weq{16}{N_a(\varphi) = a_{1,1}a_{2,2},}
und $a_{1,1}$ ist die kleinste positive ganze rationale Zahl, die durch $\varphi$ theilbar ist.

Nun k\"onnen wir andererseits zeigen, dass diese Zahl $= ec$ ist, woraus $ec = a_{1,1}$, und nach (14) und (16) $e = a_{2,2}$ folgt.

Wenn n\"amlich eine ganze rationale Zahl $m$ durch $\varphi$ theilbar sein soll, so muss, wenn $\omega'$ zu $\omega$ conjugirt ist, da
\[
ct_1^2 + bt_1t_2 - at_2^2
\]
ein Einheitsfunctional ist,
\[
\frac{m(ct_1^2+bt_1t_2-at_2^2)}{\varphi} = \frac{m}{e}(t_1+\omega' t_2)
\]
ein ganzes Functional sein. Folglich m\"ussen
\[
\frac{m}{e} \quad \text{und} \quad \frac{m}{e}\omega'
\]
ganze Zahlen sein, woraus folgt, dass $m : e$ eine durch $c$ theilbare, also $m$ eine durch $ec$ theilbare ganze rationale Zahl ist. Da andererseits $ec$ durch $\varphi$ theilbar ist, so folgt, dass $ec$ die kleinste diesen Forderungen gen\"ugende Zahl ist. Ausserdem ergiebt sich noch, dass die ganze Zahl $ec\omega$ durch $\varphi$ theilbar ist. Denn $ec\omega$ ist in der Darstellung (12) ein Co\"efficient der ganzen Function $\varphi$ und daher durch $\varphi$ theilbar (\S.\,142).

Dr\"ucken wir $\Theta$ in (15) durch $\omega$ aus, so ergiebt sich nach (11)
\[
\begin{aligned}
\alpha_2 &= a_{1,2} - \frac{eb}{2} + ec\omega, \quad b \text{ gerade}\\
&= a_{1,2} - \frac{e(b-1)}{2} + ec\omega, \quad b \text{ ungerade,}
\end{aligned}
\]
und es ist also
\[
a_{1,2} \equiv \frac{eb}{2} \quad \text{oder} \quad \equiv \frac{e(b-1)}{2}
\]
nach dem Modul $\varphi$, und folglich, da beide Seiten rational sind, auch nach dem Modul $a_{1,1}=ec$. Demnach kann $a_{1,2}$ dem Werthe $\frac{1}{2}eb$ oder $\frac{1}{2}e(b-1)$ gleich gesetzt werden, und wir erhalten $\alpha_2 = ec\omega$, was zu beweisen war.

Um das hierdurch Bewiesene \"ubersichtlich zusammenzufassen, k\"onnen wir folgende S\"atze aussprechen:
%CONT>>
